\section{AHSME 1958}
        \begin{enumerate}
        	\item (\textit{AHSME 1958}) The value of $[2 - 3(2 - 3)^{-1}]^{-1}$ is:


 $\textbf{(A)}\ 5\qquad  \textbf{(B)}\ -5\qquad  \textbf{(C)}\ \frac{1}{5}\qquad  \textbf{(D)}\ -\frac{1}{5}\qquad  \textbf{(E)}\ \frac{5}{3}$


	\item (\textit{AHSME 1958}) If $\frac {1}{x} - \frac {1}{y} = \frac {1}{z}$, then $z$ equals:


 $\textbf{(A)}\ y - x\qquad \textbf{(B)}\ x - y\qquad \textbf{(C)}\ \frac {y - x}{xy}\qquad \textbf{(D)}\ \frac {xy}{y - x}\qquad \textbf{(E)}\ \frac {xy}{x - y}$


	\item (\textit{AHSME 1958}) Of the following expressions the one equal to $\frac{a^{-1}b^{-1}}{a^{-3} - b^{-3}}$ is:


 $\textbf{(A)}\ \frac{a^2b^2}{b^2 - a^2}\qquad  \textbf{(B)}\ \frac{a^2b^2}{b^3 - a^3}\qquad  \textbf{(C)}\ \frac{ab}{b^3 - a^3}\qquad  \textbf{(D)}\ \frac{a^3 - b^3}{ab}\qquad  \textbf{(E)}\ \frac{a^2b^2}{a - b}$


	\item (\textit{AHSME 1958}) In the expression $\frac{x + 1}{x - 1}$ each $x$ is replaced by $\frac{x + 1}{x - 1}$. The resulting expression, evaluated for $x = \frac{1}{2}$, equals:


 $\textbf{(A)}\ 3\qquad  \textbf{(B)}\ -3\qquad  \textbf{(C)}\ 1\qquad  \textbf{(D)}\ -1\qquad  \textbf{(E)}\ \text{none of these}$


	\item (\textit{AHSME 1958}) The expression $2 + \sqrt{2} + \frac{1}{2 + \sqrt{2}} + \frac{1}{\sqrt{2} - 2}$  equals:


 $\textbf{(A)}\ 2\qquad  \textbf{(B)}\ 2 - \sqrt{2}\qquad  \textbf{(C)}\ 2 + \sqrt{2}\qquad  \textbf{(D)}\ 2\sqrt{2}\qquad  \textbf{(E)}\ \frac{\sqrt{2}}{2}$


	\item (\textit{AHSME 1958}) The arithmetic mean between $\frac {x + a}{x}$ and $\frac {x - a}{x}$, when $x \neq 0$, is:


 $\textbf{(A)}\ {2}\text{, if }{a \neq 0}\qquad \textbf{(B)}\ 1\qquad \textbf{(C)}\ {1}\text{, only if }{a = 0}\qquad \textbf{(D)}\ \frac {a}{x}\qquad \textbf{(E)}\ x$


	\item (\textit{AHSME 1958}) A straight line joins the points $(-1,1)$ and $(3,9)$. Its $x$-intercept is:


 $\textbf{(A)}\ -\frac{3}{2}\qquad  \textbf{(B)}\ -\frac{2}{3}\qquad  \textbf{(C)}\ \frac{2}{5}\qquad  \textbf{(D)}\ 2\qquad  \textbf{(E)}\ 3$


	\item (\textit{AHSME 1958}) Which of these four numbers $\sqrt{\pi^2},\,\sqrt[3]{.8},\,\sqrt[4]{.00016},\,\sqrt[3]{-1}\cdot \sqrt{(.09)^{-1}}$, is (are) rational:


 $\textbf{(A)}\ \text{none}\qquad  \textbf{(B)}\ \text{all}\qquad  \textbf{(C)}\ \text{the first and fourth}\qquad  \textbf{(D)}\ \text{only the fourth}\qquad  \textbf{(E)}\ \text{only the first}$


	\item (\textit{AHSME 1958}) A value of $x$ satisfying the equation $x^2 + b^2 = (a - x)^2$ is:


 $\textbf{(A)}\ \frac{b^2 + a^2}{2a}\qquad  \textbf{(B)}\ \frac{b^2 - a^2}{2a}\qquad  \textbf{(C)}\ \frac{a^2 - b^2}{2a}\qquad  \textbf{(D)}\ \frac{a - b}{2}\qquad  \textbf{(E)}\ \frac{a^2 - b^2}{2}$


	\item (\textit{AHSME 1958}) For what real values of $k$, other than $k = 0$, does the equation $x^2 + kx + k^2 = 0$ have real roots?


 $\textbf{(A)}\ {k < 0}\qquad  \textbf{(B)}\ {k > 0} \qquad  \textbf{(C)}\ {k \ge 1} \qquad  \textbf{(D)}\ \text{all values of }{k}\qquad  \textbf{(E)}\ \text{no values of }{k}$


	\item (\textit{AHSME 1958}) The number of roots satisfying the equation $\sqrt{5 - x} = x\sqrt{5 - x}$ is:


 $\textbf{(A)}\ \text{unlimited}\qquad  \textbf{(B)}\ 3\qquad  \textbf{(C)}\ 2\qquad  \textbf{(D)}\ 1\qquad  \textbf{(E)}\ 0$


	\item (\textit{AHSME 1958}) If $P = \frac{s}{(1 + k)^n}$ then $n$ equals:


 $\textbf{(A)}\ \frac{\log{\left(\frac{s}{P}\right)}}{\log{(1 + k)}}\qquad  \textbf{(B)}\ \log{\left(\frac{s}{P(1 + k)}\right)}\qquad  \textbf{(C)}\ \log{\left(\frac{s - P}{1 + k}\right)}\qquad \\ \textbf{(D)}\ \log{\left(\frac{s}{P}\right)} + \log{(1 + k)}\qquad  \textbf{(E)}\ \frac{\log{(s)}}{\log{(P(1 + k))}}$


	\item (\textit{AHSME 1958}) The sum of two numbers is $10$; their product is $20$. The sum of their reciprocals is:


 $\textbf{(A)}\ \frac{1}{10}\qquad  \textbf{(B)}\ \frac{1}{2}\qquad  \textbf{(C)}\ 1\qquad  \textbf{(D)}\ 2\qquad  \textbf{(E)}\ 4$


	\item (\textit{AHSME 1958}) At a dance party a group of boys and girls exchange dances as follows: one boy dances with $5$ girls, a second boy dances with $6$ girls, and so on, the last boy dancing with all the girls. If $b$ represents the number of boys and $g$ the number of girls, then:


 $\textbf{(A)}\ b = g\qquad  \textbf{(B)}\ b = \frac{g}{5}\qquad  \textbf{(C)}\ b = g - 4\qquad  \textbf{(D)}\ b = g - 5\qquad \\ \textbf{(E)}\ \text{It is impossible to determine a relation between }{b}\text{ and }{g}\text{ without knowing }{b + g.}$


	\item (\textit{AHSME 1958}) A quadrilateral is inscribed in a circle. If an angle is inscribed into each of the four segments outside the quadrilateral, the sum of these four angles, expressed in degrees, is:


 $\textbf{(A)}\ 1080\qquad  \textbf{(B)}\ 900\qquad  \textbf{(C)}\ 720\qquad  \textbf{(D)}\ 540\qquad  \textbf{(E)}\ 360$


	\item (\textit{AHSME 1958}) The area of a circle inscribed in a regular hexagon is $100\pi$. The area of hexagon is:


 $\textbf{(A)}\ 600\qquad  \textbf{(B)}\ 300\qquad  \textbf{(C)}\ 200\sqrt{2}\qquad  \textbf{(D)}\ 200\sqrt{3}\qquad  \textbf{(E)}\ 120\sqrt{5}$


	\item (\textit{AHSME 1958}) If $x$ is positive and $\log{x} \ge \log{2} + \frac{1}{2}\log{x}$, then:


 $\textbf{(A)}\ {x}\text{ has no minimum or maximum value}\qquad \\ \textbf{(B)}\ \text{the maximum value of }{x}\text{ is }{1}\qquad \\ \textbf{(C)}\ \text{the minimum value of }{x}\text{ is }{1}\qquad \\ \textbf{(D)}\ \text{the maximum value of }{x}\text{ is }{4}\qquad \\ \textbf{(E)}\ \text{the minimum value of }{x}\text{ is }{4}$


	\item (\textit{AHSME 1958}) The area of a circle is doubled when its radius $r$ is increased by $n$. Then $r$ equals:


 $\textbf{(A)}\ n(\sqrt{2} + 1)\qquad  \textbf{(B)}\ n(\sqrt{2} - 1)\qquad  \textbf{(C)}\ n\qquad  \textbf{(D)}\ n(2 - \sqrt{2})\qquad  \textbf{(E)}\ \frac{n\pi}{\sqrt{2} + 1}$


	\item (\textit{AHSME 1958}) The sides of a right triangle are $a$ and $b$ and the hypotenuse is $c$. A perpendicular from the vertex divides $c$ into segments $r$ and $s$, adjacent respectively to $a$ and $b$. If $a : b = 1 : 3$, then the ratio of $r$ to $s$ is:


 $\textbf{(A)}\ 1 : 3\qquad  \textbf{(B)}\ 1 : 9\qquad  \textbf{(C)}\ 1 : 10\qquad  \textbf{(D)}\ 3 : 10\qquad  \textbf{(E)}\ 1 : \sqrt{10}$


	\item (\textit{AHSME 1958}) If $4^x - 4^{x - 1} = 24$, then $(2x)^x$  equals:


 $\textbf{(A)}\ 5\sqrt{5}\qquad  \textbf{(B)}\ \sqrt{5}\qquad  \textbf{(C)}\ 25\sqrt{5}\qquad  \textbf{(D)}\ 125\qquad  \textbf{(E)}\ 25$


	\item (\textit{AHSME 1958}) In the accompanying figure $\overline{CE}$ and $\overline{DE}$ are equal chords of a circle with center $O$. Arc $AB$ is a quarter-circle. Then the ratio of the area of triangle $CED$ to the area of triangle $AOB$ is:


 
                    \begin{figure}[H]
                        \centering
                        \includegraphics[width=0.5\textwidth]{Problem/AHSME_I_1958_21_Problem_Figure_1.png}
                    \end{figure}
                    
$\textbf{(A)}\ \sqrt {2} : 1\qquad \textbf{(B)}\ \sqrt {3} : 1\qquad \textbf{(C)}\ 4 : 1\qquad \textbf{(D)}\ 3 : 1\qquad \textbf{(E)}\ 2 : 1$


	\item (\textit{AHSME 1958}) A particle is placed on the parabola $y = x^2 - x - 6$ at a point $P$ whose $y$-coordinate is $6$. It is allowed to roll along the parabola until it reaches the nearest point $Q$ whose $y$-coordinate is $-6$. The horizontal distance traveled by the particle (the numerical value of the difference in the $x$-coordinates of $P$ and $Q$) is:


 $\textbf{(A)}\ 5\qquad  \textbf{(B)}\ 4\qquad  \textbf{(C)}\ 3\qquad  \textbf{(D)}\ 2\qquad  \textbf{(E)}\ 1$


	\item (\textit{AHSME 1958}) If, in the expression $x^2 - 3$, $x$ increases or decreases by a positive amount of $a$, the expression changes by an amount:


 $\textbf{(A)}\ {\pm 2ax + a^2}\qquad  \textbf{(B)}\ {2ax \pm a^2}\qquad  \textbf{(C)}\ {\pm a^2 - 3} \qquad  \textbf{(D)}\ {(x + a)^2 - 3}\qquad\\  \textbf{(E)}\ {(x - a)^2 - 3}$


	\item (\textit{AHSME 1958}) A man travels $m$ feet due north at $2$ minutes per mile. He returns due south to his starting point at $2$ miles per minute. The average rate in miles per hour for the entire trip is:


 $\textbf{(A)}\ 75\qquad  \textbf{(B)}\ 48\qquad  \textbf{(C)}\ 45\qquad  \textbf{(D)}\ 24\qquad\\  \textbf{(E)}\ \text{impossible to determine without knowing the value of }{m}$


	\item (\textit{AHSME 1958}) If $\log_{k}{x}\cdot \log_{5}{k} = 3$, then $x$ equals:


 $\textbf{(A)}\ k^6\qquad  \textbf{(B)}\ 5k^3\qquad  \textbf{(C)}\ k^3\qquad  \textbf{(D)}\ 243\qquad  \textbf{(E)}\ 125$


	\item (\textit{AHSME 1958}) A set of $n$ numbers has the sum $s$. Each number of the set is increased by $20$, then multiplied by $5$, and then decreased by $20$. The sum of the numbers in the new set thus obtained is:


 $\textbf{(A)}\ s + 20n\qquad  \textbf{(B)}\ 5s + 80n\qquad  \textbf{(C)}\ s\qquad  \textbf{(D)}\ 5s\qquad  \textbf{(E)}\ 5s + 4n$


	\item (\textit{AHSME 1958}) The points $(2,-3)$, $(4,3)$, and $(5, k/2)$ are on the same straight line. The value(s) of $k$ is (are):


 $\textbf{(A)}\ 12\qquad  \textbf{(B)}\ -12\qquad  \textbf{(C)}\ \pm 12\qquad  \textbf{(D)}\ {12}\text{ or }{6}\qquad  \textbf{(E)}\ {6}\text{ or }{6\frac{2}{3}}$


	\item (\textit{AHSME 1958}) A $16$-quart radiator is filled with water. Four quarts are removed and replaced with pure antifreeze liquid. Then four quarts of the mixture are removed and replaced with pure antifreeze. This is done a third and a fourth time. The fractional part of the final mixture that is water is:


 $\textbf{(A)}\ \frac{1}{4}\qquad  \textbf{(B)}\ \frac{81}{256}\qquad  \textbf{(C)}\ \frac{27}{64}\qquad  \textbf{(D)}\ \frac{37}{64}\qquad  \textbf{(E)}\ \frac{175}{256}$


	\item (\textit{AHSME 1958}) In a general triangle $ADE$ (as shown) lines $\overline{EB}$ and $\overline{EC}$ are drawn. Which of the following angle relations is true?


 
                    \begin{figure}[H]
                        \centering
                        \includegraphics[width=0.5\textwidth]{Problem/AHSME_I_1958_29_Problem_Figure_1.png}
                    \end{figure}
                    
$\textbf{(A)}\ x + z = a + b\qquad \textbf{(B)}\ y + z = a + b\qquad \textbf{(C)}\ m + x = w + n\qquad \\ \textbf{(D)}\ x + z + n = w + c + m\qquad \textbf{(E)}\ x + y + n = a + b + m$


	\item (\textit{AHSME 1958}) If $xy = b$ and $\frac{1}{x^2} + \frac{1}{y^2} = a$, then $(x + y)^2$ equals:


 $\textbf{(A)}\ (a + 2b)^2\qquad  \textbf{(B)}\ a^2 + b^2\qquad  \textbf{(C)}\ b(ab + 2)\qquad  \textbf{(D)}\ ab(b + 2)\qquad  \textbf{(E)}\ \frac{1}{a} + 2b$


	\item (\textit{AHSME 1958}) The altitude drawn to the base of an isosceles triangle is $8$, and the perimeter $32$. The area of the triangle is:


 $\textbf{(A)}\ 56\qquad  \textbf{(B)}\ 48\qquad  \textbf{(C)}\ 40\qquad  \textbf{(D)}\ 32\qquad  \textbf{(E)}\ 24$


	\item (\textit{AHSME 1958}) With $ $1000$ a rancher is to buy steers at $ $25$ each and cows at $ $26$ each. If the number of steers $s$ and the number of cows $c$ are both positive integers, then:


 $\textbf{(A)}\ \text{this problem has no solution}\qquad\\  \textbf{(B)}\ \text{there are two solutions with }{s}\text{ exceeding }{c}\qquad \\ \textbf{(C)}\ \text{there are two solutions with }{c}\text{ exceeding }{s}\qquad \\ \textbf{(D)}\ \text{there is one solution with }{s}\text{ exceeding }{c}\qquad \\ \textbf{(E)}\ \text{there is one solution with }{c}\text{ exceeding }{s}$


	\item (\textit{AHSME 1958}) For one root of $ax^2 + bx + c = 0$ to be double the other, the coefficients $a,\,b,\,c$ must be related as follows:


 $\textbf{(A)}\ 4b^2 = 9c\qquad  \textbf{(B)}\ 2b^2 = 9ac\qquad  \textbf{(C)}\ 2b^2 = 9a\qquad \\ \textbf{(D)}\ b^2 - 8ac = 0\qquad  \textbf{(E)}\ 9b^2 = 2ac$


	\item (\textit{AHSME 1958}) The numerator of a fraction is $6x + 1$, then denominator is $7 - 4x$, and $x$ can have any value between $-2$ and $2$, both included. The values of $x$ for which the numerator is greater than the denominator are:


 $\textbf{(A)}\ \frac{3}{5} < x \le 2\qquad  \textbf{(B)}\ \frac{3}{5} \le x \le 2\qquad  \textbf{(C)}\ 0 < x \le 2\qquad \\ \textbf{(D)}\ 0 \le x \le 2\qquad  \textbf{(E)}\ -2 \le x \le 2$


	\item (\textit{AHSME 1958}) A triangle is formed by joining three points whose coordinates are integers. If the $x$-coordinate and the $y$-coordinate each have a value of $1$, then the area of the triangle, in square units:


 $\textbf{(A)}\ \text{must be an integer}\qquad  \textbf{(B)}\ \text{may be irrational}\qquad  \textbf{(C)}\ \text{must be irrational}\qquad  \textbf{(D)}\ \text{must be rational}\qquad \\ \textbf{(E)}\ \text{will be an integer only if the triangle is equilateral.}$


	\item (\textit{AHSME 1958}) The sides of a triangle are $30$, $70$, and $80$ units. If an altitude is dropped upon the side of length $80$, the larger segment cut off on this side is:


 $\textbf{(A)}\ 62\qquad  \textbf{(B)}\ 63\qquad  \textbf{(C)}\ 64\qquad  \textbf{(D)}\ 65\qquad  \textbf{(E)}\ 66$


	\item (\textit{AHSME 1958}) The first term of an arithmetic series of consecutive integers is $k^2 + 1$. The sum of $2k + 1$ terms of this series may be expressed as:


 $\textbf{(A)}\ k^3 + (k + 1)^3\qquad  \textbf{(B)}\ (k - 1)^3 + k^3\qquad  \textbf{(C)}\ (k + 1)^3\qquad \\ \textbf{(D)}\ (k + 1)^2\qquad  \textbf{(E)}\ (2k + 1)(k + 1)^2$


	\item (\textit{AHSME 1958}) Let $r$ be the distance from the origin to a point $P$ with coordinates $x$ and $y$. Designate the ratio $\frac{y}{r}$ by $s$ and the ratio $\frac{x}{r}$  by $c$. Then the values of $s^2 - c^2$ are limited to the numbers:


 $\textbf{(A)}\ \text{less than }{-1}\text{ are greater than }{+1}\text{, both excluded}\qquad\\ \textbf{(B)}\ \text{less than }{-1}\text{ are greater than }{+1}\text{, both included}\qquad \\ \textbf{(C)}\ \text{between }{-1}\text{ and }{+1}\text{, both excluded}\qquad \\ \textbf{(D)}\ \text{between }{-1}\text{ and }{+1}\text{, both included}\qquad \\ \textbf{(E)}\ {-1}\text{ and }{+1}\text{ only}$


	\item (\textit{AHSME 1958}) We may say concerning the solution of $|x|^2 + |x| - 6 =0$ that:


 $\textbf{(A)}\ \text{there is only one root}\qquad  \textbf{(B)}\ \text{the sum of the roots is }{+1}\qquad  \textbf{(C)}\ \text{the sum of the roots is }{0}\qquad \\ \textbf{(D)}\ \text{the product of the roots is }{+4}\qquad  \textbf{(E)}\ \text{the product of the roots is }{-6}$


	\item (\textit{AHSME 1958}) Given $a_0 = 1$, $a_1 = 3$, and the general relation $a_n^2 - a_{n - 1}a_{n + 1} = (-1)^n$ for $n \ge 1$. Then $a_3$ equals:


 $\textbf{(A)}\ \frac{13}{27}\qquad  \textbf{(B)}\ 33\qquad  \textbf{(C)}\ 21\qquad  \textbf{(D)}\ 10\qquad  \textbf{(E)}\ -17$


	\item (\textit{AHSME 1958}) The roots of $Ax^2 + Bx + C = 0$ are $r$ and $s$. For the roots of
$x^2+px +q =0$


 to be $r^2$ and $s^2$, $p$ must equal:


 $\textbf{(A)}\ \frac{B^2 - 4AC}{A^2}\qquad  \textbf{(B)}\ \frac{B^2 - 2AC}{A^2}\qquad  \textbf{(C)}\ \frac{2AC - B^2}{A^2}\qquad \\ \vspace{0.1in}\\ \textbf{(D)}\ B^2 - 2C\qquad  \textbf{(E)}\ 2C - B^2$


	\item (\textit{AHSME 1958}) In a circle with center $O$, chord $\overline{AB}$ equals chord $\overline{AC}$. Chord $\overline{AD}$ cuts $\overline{BC}$ in $E$. If $AC = 12$ and $AE = 8$, then $AD$ equals:


 $\textbf{(A)}\ 27\qquad  \textbf{(B)}\ 24\qquad  \textbf{(C)}\ 21\qquad  \textbf{(D)}\ 20\qquad  \textbf{(E)}\ 18$


	\item (\textit{AHSME 1958}) $\overline{AB}$ is the hypotenuse of a right triangle $ABC$. Median $\overline{AD}$ has length $7$ and median $\overline{BE}$ has length $4$. The length of $\overline{AB}$ is:


 $\textbf{(A)}\ 10\qquad  \textbf{(B)}\ 5\sqrt{3}\qquad  \textbf{(C)}\ 5\sqrt{2}\qquad  \textbf{(D)}\ 2\sqrt{13}\qquad  \textbf{(E)}\ 2\sqrt{15}$


	\item (\textit{AHSME 1958}) Given the true statements: (1) If $a$ is greater than $b$, then $c$ is greater than $d$ (2) If $c$ is less than $d$, then $e$ is greater than $f$. A valid conclusion is:


 $\textbf{(A)}\ \text{If }{a}\text{ is less than }{b}\text{, then }{e}\text{ is greater than }{f}\qquad \\ \textbf{(B)}\ \text{If }{e}\text{ is greater than }{f}\text{, then }{a}\text{ is less than }{b}\qquad \\ \textbf{(C)}\ \text{If }{e}\text{ is less than }{f}\text{, then }{a}\text{ is greater than }{b}\qquad \\ \textbf{(D)}\ \text{If }{a}\text{ is greater than }{b}\text{, then }{e}\text{ is less than }{f}\qquad \\ \textbf{(E)}\ \text{none of these}$


	\item (\textit{AHSME 1958}) A check is written for $x$ dollars and $y$ cents, $x$ and $y$ both two-digit numbers. In error it is cashed for $y$ dollars and $x$ cents, the incorrect amount exceeding the correct amount by $\$17.82$. Then:


 $\textbf{(A)}\ {x}\text{ cannot exceed }{70}\qquad \\ \textbf{(B)}\ {y}\text{ can equal }{2x}\qquad\\ \textbf{(C)}\ \text{the amount of the check cannot be a multiple of }{5}\qquad \\ \textbf{(D)}\ \text{the incorrect amount can equal twice the correct amount}\qquad \\ \textbf{(E)}\ \text{the sum of the digits of the correct amount is divisible by }{9}$


	\item (\textit{AHSME 1958}) For values of $x$ less than $1$ but greater than $-4$, the expression
$\frac{x^2 - 2x + 2}{2x - 2}$
has:


 $\textbf{(A)}\ \text{no maximum or minimum value}\qquad \\ \textbf{(B)}\ \text{a minimum value of }{+1}\qquad \\ \textbf{(C)}\ \text{a maximum value of }{+1}\qquad \\ \textbf{(D)}\ \text{a minimum value of }{-1}\qquad \\ \textbf{(E)}\ \text{a maximum value of }{-1}$


	\item (\textit{AHSME 1958}) $ABCD$ is a rectangle (see the accompanying diagram) with $P$ any point on $\overline{AB}$. $\overline{PS} \perp \overline{BD}$ and $\overline{PR} \perp \overline{AC}$. $\overline{AF} \perp \overline{BD}$ and $\overline{PQ} \perp \overline{AF}$. Then $PR + PS$ is equal to:


 
                    \begin{figure}[H]
                        \centering
                        \includegraphics[width=0.5\textwidth]{Problem/AHSME_I_1958_47_Problem_Figure_1.png}
                    \end{figure}
                    
$\textbf{(A)}\ PQ\qquad  \textbf{(B)}\ AE\qquad  \textbf{(C)}\ PT + AT\qquad  \textbf{(D)}\ AF\qquad  \textbf{(E)}\ EF$


	\item (\textit{AHSME 1958}) Diameter $\overline{AB}$ of a circle with center $O$ is $10$ units. $C$ is a point $4$ units from $A$, and on $\overline{AB}$. $D$ is a point $4$ units from $B$, and on $\overline{AB}$. $P$ is any point on the circle. Then the broken-line path from $C$ to $P$ to $D$:


 $\textbf{(A)}\ \text{has the same length for all positions of }{P}\qquad\\  \textbf{(B)}\ \text{exceeds }{10}\text{ units for all positions of }{P}\qquad \\ \textbf{(C)}\ \text{cannot exceed }{10}\text{ units}\qquad \\ \textbf{(D)}\ \text{is shortest when }{\triangle CPD}\text{ is a right triangle}\qquad \\ \textbf{(E)}\ \text{is longest when }{P}\text{ is equidistant from }{C}\text{ and }{D}.$


	\item (\textit{AHSME 1958}) In the expansion of $(a + b)^n$ there are $n + 1$ dissimilar terms. The number of dissimilar terms in the expansion of $(a + b + c)^{10}$ is:


 $\textbf{(A)}\ 11\qquad  \textbf{(B)}\ 33\qquad  \textbf{(C)}\ 55\qquad  \textbf{(D)}\ 66\qquad  \textbf{(E)}\ 132$


	\item (\textit{AHSME 1958}) In this diagram a scheme is indicated for associating all the points of segment $\overline{AB}$ with those of segment $\overline{A'B'}$, and reciprocally. To described this association scheme analytically, let $x$ be the distance from a point $P$ on $\overline{AB}$ to $D$ and let $y$ be the distance from the associated point $P'$ of $\overline{A'B'}$ to $D'$. Then for any pair of associated points, if $x = a,\, x + y$ equals:


 
                    \begin{figure}[H]
                        \centering
                        \includegraphics[width=0.5\textwidth]{Problem/AHSME_I_1958_50_Problem_Figure_1.png}
                    \end{figure}
                    
$\textbf{(A)}\ 13a\qquad \textbf{(B)}\ 17a - 51\qquad \textbf{(C)}\ 17 - 3a\qquad \textbf{(D)}\ \frac {17 - 3a}{4}\qquad \textbf{(E)}\ 12a - 34$


\end{enumerate}